\chapter{Endometritis}
\index{Endometritis}

\section*{Definition and Overview}
Endometritis is inflammation or infection of the lining of the uterus, called the endometrium. It most commonly occurs after childbirth, miscarriage, abortion, or gynaecological procedures, when bacteria enter the uterus. It can also be caused by sexually transmitted infections and, in a chronic form, is associated with pelvic inflammatory disease and fertility problems. Acute endometritis is usually treated successfully with antibiotics, and prompt treatment prevents serious complications such as pelvic abscess and sepsis.

\section*{Epidemiology}
Postpartum endometritis is the most common form and is an important cause of fever after childbirth, particularly after caesarean section. The risk is higher after prolonged labour, prolonged rupture of membranes, and retained placental tissue. Endometritis after miscarriage, abortion, or procedures such as hysteroscopy is less common but well recognised. Chronic endometritis is increasingly recognised as a cause of recurrent miscarriage and implantation failure.

\section*{Aetiology and Pathophysiology}
Endometritis occurs when bacteria enter the normally sterile uterus.
\begin{itemize}
    \item \textbf{Postpartum:} bacteria ascend from the vagina after childbirth, especially after caesarean, retained membranes, or long labour
    \item \textbf{Post-procedure:} after miscarriage management, abortion, hysteroscopy, or insertion of an intrauterine device
    \item \textbf{STIs:} chlamydia and gonorrhoea can cause infection of the uterine lining
    \item \textbf{Organisms:} a mixture of vaginal and bowel bacteria, including streptococci, staphylococci, and anaerobes
    \item \textbf{Mechanism:} bacteria infect the endometrium, causing inflammation, which can spread to the muscle of the uterus, the tubes, and the pelvis
    \item \textbf{Chronic endometritis:} persistent inflammation, often with no clear acute illness, linked to recurrent pregnancy loss
\end{itemize}

\section*{Clinical Features}
Symptoms usually develop within days of the triggering event.
\begin{itemize}
    \item Fever and chills
    \item Lower abdominal or pelvic pain
    \item Abnormal or foul-smelling vaginal discharge
    \item Heavy or prolonged bleeding, including after childbirth
    \item A tender uterus on examination
    \item General malaise and feeling unwell
    \item Chronic endometritis may cause no symptoms or only mild, ongoing discomfort
\end{itemize}

\section*{Differential Diagnosis}
Other causes of pelvic infection and pain must be considered.
\begin{itemize}
    \item Pelvic inflammatory disease affecting the tubes and ovaries
    \item Urinary tract infection and kidney infection
    \item Wound infection after caesarean section
    \item Mastitis in breastfeeding women
    \item Deep vein thrombosis after childbirth
    \item Retained products of conception
    \item Other causes of post-partum fever, including breast infection and chest infection
\end{itemize}

\section*{When to Seek Medical Care}
Any fever, pelvic pain, or abnormal discharge after childbirth, miscarriage, abortion, or a gynaecological procedure should be assessed promptly by a doctor.

\textbf{Contact your local emergency services for:}
\begin{itemize}
    \item High fever with severe pelvic pain
    \item Heavy bleeding with faintness or dizziness
    \item Foul-smelling discharge with signs of being unwell
    \item Confusion, collapse, or signs of sepsis
\end{itemize}

\section*{Diagnosis and Investigations}
Diagnosis is usually clinical, with tests to confirm and exclude other causes.
\begin{itemize}
    \item \textbf{History:} recent childbirth, miscarriage, or procedures
    \item \textbf{Examination:} fever, pelvic tenderness, and discharge
    \item \textbf{Blood tests:} full blood count and inflammatory markers
    \item \textbf{Swabs:} cervical and vaginal swabs for infection
    \item \textbf{Blood cultures:} in severe illness
    \item \textbf{Ultrasound:} to check for retained tissue or abscess
    \item \textbf{Endometrial biopsy:} for chronic endometritis, to confirm inflammation
\end{itemize}

\section*{Management}
Treatment is with antibiotics, and any retained tissue is removed.
\begin{itemize}
    \item \textbf{Antibiotics:} broad-spectrum antibiotics, given intravenously in hospital for moderate to severe infection
    \item \textbf{Removal of retained tissue:} surgical evacuation of retained products of conception
    \item \textbf{Pain relief:} anti-inflammatory medicines
    \item \textbf{Rest and fluids:} supportive care during recovery
    \item \textbf{Treatment of STIs:} antibiotics for chlamydia or gonorrhoea, with partner treatment
    \item \textbf{Chronic endometritis:} a course of antibiotics, which may improve fertility outcomes
\end{itemize}

\section*{Complications}
\begin{itemize}
    \item Spread of infection to the pelvic organs, causing pelvic inflammatory disease
    \item Pelvic abscess
    \item Sepsis, a life-threatening whole-body infection
    \item Chronic pelvic pain
    \item Scarring and adhesions, affecting future fertility
    \item Recurrent miscarriage in chronic endometritis
\end{itemize}

\section*{Prognosis and Outcomes}
With prompt antibiotic treatment, most cases of endometritis resolve fully within a week or two, and serious complications are uncommon. Chronic endometritis responds to antibiotics in most women, and fertility outcomes can improve after treatment. Early recognition and treatment are the keys to a good outcome.

\section*{Prevention and Self-Care}
\begin{itemize}
    \item Practise safe sex and get tested for sexually transmitted infections
    \item Attend postnatal checks and seek care for any fever
    \item Follow advice after miscarriage, abortion, and gynaecological procedures
    \item Maintain good hygiene and use pads rather than tampons while bleeding
    \item Complete the full course of antibiotics if prescribed
    \item Seek prompt care for any signs of infection
\end{itemize}

\section*{Sources and Further Reading}
The following reputable sources provide further information for healthcare students and patients seeking reliable, current advice about endometritis.
\begin{itemize}
    \item The Royal Women's Hospital. \textit{Postpartum infection information.} https://www.thewomens.org.au/
    \item Healthdirect Australia. \textit{Pelvic inflammatory disease.} https://www.healthdirect.gov.au/pelvic-inflammatory-disease
    \item Jean Hailes for Women's Health. \textit{Pelvic infection information.} https://www.jeanhailes.org.au/
    \item NHS. \textit{Pelvic inflammatory disease.} https://www.nhs.uk/conditions/pelvic-inflammatory-disease/
    \item Mayo Clinic. \textit{Endometritis and pelvic infections.} https://www.mayoclinic.org/diseases-conditions/
    \item MSD Manual. \textit{Endometritis: professional overview.} https://www.msdmanuals.com/
\end{itemize}
